\chapter{LEEP (Loop Electrosurgical Excision Procedure)}

\section*{Overview}
LEEP removes abnormal cervical tissue using a thin wire loop carrying an electrical current. The tissue is sent to a laboratory to diagnose and treat precancerous cervical changes. The exact approach, setting, and preparation depend on the indication, anatomy, and the patient's health.

\section*{Why It Is Done}
It is commonly recommended for biopsy-confirmed high-grade cervical dysplasia or selected persistent abnormal cervical changes. It can both remove the abnormal area and provide a specimen for examination.

\section*{How It Is Performed}
The patient lies in the position used for a pelvic examination. After a speculum is placed, the cervix is numbed with local anesthetic. The clinician uses the energized loop to remove transformation-zone tissue, then controls bleeding with cautery or a topical agent.

\section*{Preparation}
Pregnancy should be excluded, and the clinician should be told about blood-thinning medicines or bleeding disorders. The procedure is usually scheduled when the patient is not menstruating and is generally performed in an office or outpatient setting.

\section*{Risks and Recovery}
Cramping, light bleeding, and watery or dark discharge are common for a short time. Less common risks include infection, heavier bleeding, cervical narrowing, and a small increase in the risk of preterm birth in a future pregnancy. Patients are usually advised to avoid vaginal intercourse, tampons, and douching while the cervix heals.

\section*{Outcome and Follow-up}
The pathology report describes the grade of cervical intraepithelial neoplasia and whether the margins are clear. Follow-up usually includes HPV testing and/or cervical cytology at an interval recommended by the clinician; a positive margin or persistent abnormality may require additional evaluation.

\section*{When to Seek Medical Care}
Follow the procedure team's specific instructions. Seek urgent medical assessment for severe or worsening pain, uncontrolled bleeding, fever or signs of infection, trouble breathing, chest pain, fainting, new weakness or confusion, inability to urinate, or any rapidly worsening symptom. The appropriate warning signs depend on the procedure and the patient's medical condition.

\section*{References}
\begin{enumerate}
  \item MedlinePlus. Loop electrosurgical excision procedure (LEEP). U.S. National Library of Medicine; 2025.
  \item Current Medical Diagnosis and Treatment. McGraw-Hill; 2025.
\end{enumerate}

\clearpage
